\section{A Type-Theoretic Approach}\label{sec:type-theory}

If we view descriptors as record types (which they clearly are), we can borrow
concepts from Programming Languages theory.

\subsection{Pollux as a Logical Relation}

Pollux's goal is to be able to look at a pair of data descriptors and
characterize what happens to values under the first descriptor when parsed under
the updated one. We can split this cleanly into two parts, a syntactic property
of ``compatible descriptors'' that operates solely on the descriptors and a
semantic property that values under the original descriptor are meaningful and
interpretable when parsed under the updated one. More specifically, we want to
prove that all descriptor pairs with the syntactic property have a
transformation which the semantic property. By definition, that is a logical
relation.

Pollux then becomes a type system which can characterize functions between
messages as a compatible transformation.

\[ \vDash f : d_1 \preceq d_2 \rightarrow \forall\ m \checkmark^{d_1}, bs.\ \mathrm{Encodes}_{d_1}(m, bs) \rightarrow
  \mathrm{parse}_{d_2}(bs) = f(m) \]

Then we end up with a syntactic typing judgment on the descriptors themselves,
related to what we can convert between, and a semantic typing judgment on the
transformation between the descriptors. Pulling in the recursive type view of
the descriptors, we'd end up with a typing rule something like this (likely
\emph{not} actually correct):
\[ \infer{ X : \mu\ F \preceq \mu\ G \vdash f : F(\mu\ F) \preceq G(\mu\ G) }{\vDash \mu\ X.f : \mu\ F \preceq \mu\ G} \]
In the structure of $f$, we expect that the bound variable $X$ will be used
exactly at where $f$ makes a recursive call, and exists after restructuring a
message so the lean kernel clearly sees the syntactic well-foundedness of the
recursive call. This also aligns with logical relation step-indexing, were each
index is the size of the current message. Since this is also structural, both
techniques would work.

This framing is complicated slightly by what happens when fields can be dropped,
which is already implemented in Section~\ref{sec:proto-transform}. We naturally
expect transitivity from descriptor updates, but consider this sequence of
transformations.

\begin{minted}{proto}
  message Msg {
    int32 field = 1;
    int32 other = 2;
  } 

  message Msg' {
    int32 other = 2;
  }

  message Msg'' {
    string field = 1;
    int32 other = 2;
  }
\end{minted}

Now we have a sequence of two valid updates, dropping a field and then adding a
new field, but it doesn't compose correctly since we can't move directly from
\texttt{Msg} to \texttt{Msg''} since the wire types of field 1 are different.
Fortunately, Protobuf has a mechanism for this, reserved field numbers, which
would disallow this and restore our ability to compose updates.

\subsection{Adding Recursion}

Consider the addition of recursive descriptors to the Pollux formalization. With
a descriptor rule like this one
\[
  \infer[D-Update]{ \dom{d} \subseteq \dom{d'} \\ \forall\ k \in
    \dom{d} \cap \dom{d'}.\ d[k] \propto d'[k] }{ d \ll d' }
\]
(Note that for mandatory messages, which aren't supported in Protobuf, we'd need
that $\dom{d_{mand}} \supseteq \dom{d_{mand}}$, i.e.\ that mandatory fields can
become optional.) the descriptor relation would become co-inductive since an
update to this descriptor would would require showing $d_1[k] \propto d_2[k]$ if
$k$ is the field were we see the recursive type and the field relation would
include:
\[ \infer[F-Msg]{ d_1[k] \ll d_2[k] }{ \langle k, d_1[k] \rangle \propto \langle k, d_2[k] \rangle } \]
So a derivation that $d_1[k] \ll d_2[k]$ would need to include a derivation of
itself! The natural solution here is to make $\ll$ co-inductive. 

However, if we more closely model the descriptors as proper recursive types, it
might be possible to avoid that. Consider this Protobuf descriptor:
\begin{minted}{proto}
message Stream {
  int32 hd = 1;
  Stream tl = 2;
}
\end{minted}
Note first that while it's called \texttt{Stream}, since all Protobuf fields can
be omitted and message fields always have an explicit presence indicator, we
actually have a \texttt{List} definition rather than a true co-inductive stream.

From the prospective of algebraic type theory, this type is
$F\ X = \ints \times \optt{X}$. Note that Protobuf always generates nested messages as
if they were optional, with a host language default (i.e. \texttt{nil} in Go)
and a some mechanism to explicitly check the presence of the nested message. If
we were to add a \texttt{bool} field to this descriptor, the algebraic type
would update to
$G\ X = \mathtt{int} \times \mathtt{bool} \times \optt{X}$.

\begin{minted}{proto}
message Stream' {
  int32 hd = 1;
  bool other = 3;
  Stream tl = 2;
}
\end{minted}

Then the descriptor relation I'm proposing would be a small proof system where
we want to show that two descriptors are compatible if $\Sigma, \varnothing \vdash d \ll d'$, but we could
\emph{define} what it means for two descriptor to be compatible as 
\begin{mathpar}
  d_1 := \mu\ X.\ F\ X 

  d_2 := \mu\ X.\ G\ X 

  d_1 \ll d_2 := \exists\ f: \denote{d_1} \rightarrow \denote{d_2}.\ \forall\ m \checkmark^{d_1}, bs.\
  \mathrm{Encodes}_{d_1}(m, bs) \rightarrow \mathrm{parse}_{d_2}(bs) = f(m)
\end{mathpar}

In English, then the definition of two descriptors being compatible is that
there exists a transformation such that for all valid messages under descriptor
$d_1$, we can serialize that message into some encoding that can be parsed to
produce the image of $m$ under the transformation. Rather than trying to prove
$\varnothing \vdash d_1 \ll d_2$, which would require co-induction, we show that 
\[ \mu\ F \ll \mu\ G \vdash F (\mu\ F) \ll G (\mu\ G) \]
which can be done with normal induction. This should be equivalent to the
original statement using the relation-based definition, when viewed through the
right lens of folding and unfolding recursive types, and is in fact a
Amadio-Cardelli style rule for subtyping equi-recursive types {\color{red}
  Citation needed}.

However, the lean formalization doesn't use the $\mu$-bound version of descriptor
types (See Section~\ref{sec:log-rel-lean}), but instead a flat symbol table
$\Sigma$. This symbol table is global and fixed by the \texttt{.proto} file,
containing names and descriptors for all the cycles the descriptor of interest
is part of. Additionally, there is $A$, a set of name pairs currently assumed to
be compatible; it grows as the derivation descends. The judgment is then $\Sigma; A
\vdash d_1 \ll d_2$, with two rules for names:

\begin{mathpar}
  \infer[N-Assum]{ (n_1, n_2) \in A }{ \Sigma; A \vdash n_1 \ll n_2 }

  \infer[N-Unfold]{ (n_1, n_2) \notin A \\
    \Sigma; A \cup \{(n_1, n_2)\} \vdash \Sigma(n_1) \ll \Sigma(n_2) }{
    \Sigma; A \vdash n_1 \ll n_2 }
\end{mathpar}

Compatibility of the top level descriptors is $\Sigma; \varnothing \vdash d_1 \ll d_2$. The
\textsc{N-Unfold} rule is equivalent to $\mu$ unfolding, with a side condition
that $A$ grows.

Co-induction is avoided for two independent finiteness facts:
\begin{enumerate}
\item Syntactically, a cyclic descriptor using a symbol table is representable
  in finite syntax. A symbol table stores \texttt{Stream}'s second field as the
  \emph{name} \texttt{Stream} rather than a copy of the descriptor. This is
  equivalent to the $\mu$-bound version. What is infinite is what the descriptor
  unfolds to, and that unfolding is a regular tree. The checker exploit this by
  carrying an assumption set $A$ of pairs of message names where
  $(n_1, n_2) \in A$ records that $\Sigma(n_1) \ll \Sigma(n_2)$ has already been assumed. It
  is what the $\varnothing$ in $\varnothing \vdash d \ll d'$ is referencing and grows
  over the course of a derivation. Before recursing on a pair of names
  encountered while checking compatibility, the checker checks if that pair is
  already in $A$ and immediately succeeds if so, adding the pair
  otherwise. Every recursive call can add a pair so $A$ never shrinks so $A$ can
  grow to at most
  $|\dom{\Sigma} \times \dom{\Sigma}| = \mathcal{O}(\dom{\Sigma}^2)$ before every recursive call would
  immediately succeed. Thus the derivation is finite.

  Thinking about this with the \texttt{Stream} example, checking \texttt{Stream}
  against \texttt{Stream'} takes two steps. In the initial check,
  $(\mathtt{Stream}, \mathtt{Stream'}) \not \in A$ so it is added and both
  descriptors are unfolded once. Field 1 gives $\ints \propto \ints$ and field 2
  reduces to $\Sigma, A \vdash \mathtt{Stream} \ll \mathtt{Stream'}$ which is now in $A$.
\item Semantically, all Protobuf values are finite even under a cyclic
  schema. Specifically, all nested messages can be omitted since they have
  presence checks and all official Protobuf implementations only traverse 100
  layers deep into a message (or 10,000 for the Go implementation for some
  reason)\footnote{https://protobuf.dev/programming-guides/proto-limits/} after
  which all parsing attempts will fail. Since our definitions always work with
  the idea that parsing will succeed, we can actually assume that the messages
  are finite and even use fuel in the proofs if needed. When the soundness proof
  re-encounters an assumed message, it has necessarily descended down at least
  one layer. Soundness can then use strong induction on the writer's message
  size. This fuel approach is also the step-indexing piece of a logical relation
  argument. We saw this with the \texttt{Stream} example, which reveals an
  important semantic fact we can take advantage of rather than a naming
  discrepancy.
\end{enumerate}

Note that an earlier definition of this featured $\exists\ bs$, which is a weaker
statement than the current version. Under the $\exists$ version, we know that there is
some way to encode $m$ so that we get $f(m)$ but for some sequence of bytes, we
won't know if a serializer was used which happens to produce that encoding. What
we really want to say is that every possible way to encode $m$ with a
spec-complaint Protobuf encoding would produce the same, known semantically
compatible value after parsing.

If the descriptor relation is a static typing judgment, then the transformation
is the semantic type judgment.

However, we even if we present this framework first, we still want to have a
relation on the descriptors for several reasons. First, the above definition
doesn't provide any insight into how to write the transformation, while we've
observed that having the relation provides a structure that the transformation
mirrors. Additionally, I personally think that the relation on descriptors
provides a nice way to check compatibility without having to construct the
transformation directly. But what is the relationship between the compatibility
relation and the transformation based definition above? The compatibility
relation is a \emph{sound} check for the semantic definition:
\[ \Sigma; \varnothing \vdash d_1 \ll d_2 \implies\ \vDash
  \transform{d_1}{\cdot}{d_2} : d_1 \preceq d_2 \]
Completeness, on the other hand, would say that every pair of descriptors
admitting \emph{some} compatible transformation is accepted by the relation. We
do not claim completeness for the logical relation, which would be harder to
prove and less useful. The compatibility relation exists to be checked, and is
intentionally conservative to rule out technically wire-compatible updates that
are undesirable like converting from \double{} to \fixeds{}. This is common with
semantic typing techniques like logical relations where the syntactic rules are
proved sound against a meaning assigned independently from them.

\subsection{Formalizing in Lean}\label{sec:log-rel-lean}

Lean doesn't have co-inductive types, unlike Rocq, and only limited support in
newer versions for co-inductive predicates and no built-in support for things
like parameterized co-induction. While the native support for co-inductive
predicates might be enough to prove what Pollux cares about, we'd be close of
the boundaries of what lean supports, and if we exceed what the built-in
co-induction mechanisms support, it quickly devolves into having to hand create
a greatest fix-point manually and state bisimulation principles manually. Rather
than that, I'd like to avoid co-induction entirely, which we've just argued is
possible.

The lean formalization could use a $\mu$-binder to actually model and interact
with recursive types, however this would be a lot of formalization and machinery
which isn't really needed. Instead, the lean formalization can use an additional
context $\Sigma$ which maps a name to a descriptor. All (mutually) recursive messages
are stored using their names in $\Sigma$ and fetched as needed. {\color{red} Would we
want all the nested messages to always be in $\Sigma$? The chains could get a lot
larger than just the 533 cycles with a maximum size of 22, but the theory should
still work out.} The compatibility relation for Protobuf
(Section~\ref{sec:proto-relations}) would need to include the $\Sigma$ context.

Additionally, the reinterpret function's termination metric should be switched
to the writer's value's size and \texttt{Desc.OneofPreservedAll} becomes a
finite conjunction over reachable pairs which mirrors the structure of the
relation itself.  % TODO

It's worth noting that since messages can always be omitted, having a cyclic
descriptor cyclic doesn't complicate Definition~\ref{def:vvalid}.

% Local Variables:
% citar-bibliography: ("../pollux.bib")
% TeX-master: "../pollux.tex"
% TeX-command-default: "Make"
% jinx-local-words: "protobuf"
% End:
